% Kandydat restrukturyzacji; nie jest aktywnym rozdziałem Atlasu.
\section{Transformacja Fouriera miar}
\providecommand{\atlasref}[2]{#2}

\subsection{Transformacja Fouriera miar}

\begin{definicja}[Transformata Fouriera skończonej miary]
\label{def:dt-fourier-measure}
Dla skończonej zespolonej miary Radona $\mu$ na $\R^d$
definiujemy
\[
 \widehat\mu(\xi)=c_d\int_{\R^d}e^{-ix\cdot\xi}
 \,\mathrm d\mu(x),\qquad c_d=(2\pi)^{-d/2}.
\]
Normą miary jest jej całkowita wariacja
$\|\mu\|=|\mu|(\R^d)$.
Por. \atlasref{ME-16::thm:complex-variation-16}{wariacja miary zespolonej}.
\end{definicja}

\begin{twierdzenie}[Ciągłość i zgodność z dystrybucjami]
\label{thm:dt-fourier-measure-cont}
Transformata $\widehat\mu$ jest ograniczona i jednostajnie
ciągła, $\|\widehat\mu\|_\infty\le c_d\|\mu\|$.
Dystrybucja temperowana
$T_\mu(\varphi)=\int\varphi\,\mathrm d\mu$
spełnia $\mathcal FT_\mu=T_{\widehat\mu}$.
Jeśli $\mu=f\lambda_d$, $f\in L^1$, to $\widehat\mu=\widehat f$.
\end{twierdzenie}
\begin{proof}
Oszacowanie wynika z nierówności
$|\int u\,\mathrm d\mu|\le\int|u|\,\mathrm d|\mu|$.
Dla dowolnego $h\in\R^d$ zachodzi
\[
 \sup_{\xi\in\R^d}|\widehat\mu(\xi+h)-\widehat\mu(\xi)|
 \le c_d\int_{\R^d}|e^{-ix\cdot h}-1|\,\mathrm d|\mu|(x).
\]
Prawa strona dąży do zera przez zbieżność zmajoryzowaną
z majorantą $2$, ponieważ $|\mu|$ jest skończona.
Ponieważ $|T_\mu(\varphi)|\le\|\mu\|\|\varphi\|_\infty$,
funkcjonał jest ciągły na $\mathcal S$.
Dla $\varphi\in\mathcal S$ moduł jądra podwójnego
parowania jest całkowalny względem
$|\mu|\otimes\lambda_d$, z całką ograniczoną przez
$c_d\|\mu\|\|\varphi\|_1$. Twierdzenie Fubiniego daje
\[
 \int_{\R^d}\widehat\mu(\xi)\varphi(\xi)\,\mathrm d\lambda_d(\xi)
 =\int_{\R^d}(\mathcal F\varphi)(x)\,\mathrm d\mu(x).
\]
To jest definicja transformaty dystrybucyjnej.
Ostatnia równość wynika z podstawienia gęstości $f$.
\end{proof}

\begin{twierdzenie}[Jednoznaczność transformaty miary]
\label{thm:dt-fourier-measure-unique}
Jeśli $\mu,\nu$ są skończonymi zespolonymi miarami Radona
i $\widehat\mu=\widehat\nu$, to $\mu=\nu$.
\end{twierdzenie}
\begin{proof}
Z Twierdzenia~\ref{thm:dt-fourier-measure-cont} mamy
$\mathcal FT_{\mu-\nu}=0$. Automorfizm $\mathcal F$ na
$\mathcal S'$ z Twierdzenia~\ref{thm:fourier-auto-tempered22}
daje $T_{\mu-\nu}=0$.
Dla $u\in C_c(\R^d)$ wybieramy nieujemną funkcję testową
$\rho$ o całce jeden i kładziemy
$\rho_\varepsilon(x)=\varepsilon^{-d}\rho(x/\varepsilon)$.
Funkcja $u*\rho_\varepsilon$ jest gładka o zwartym nośniku;
jej zbieżność jednostajna do $u$ wynika z jednostajnej
ciągłości $u$ i zwężającego się nośnika jądra:
\[
 \|u*\rho_\varepsilon-u\|_\infty
 \le\sup_{y\in\operatorname{supp}\rho_\varepsilon}
        \|\tau_yu-u\|_\infty\to0.
\]
Z ograniczoności miary otrzymujemy
$\int u\,\mathrm d(\mu-\nu)=0$ dla każdego $u\in C_c$.
Jednoznaczność w twierdzeniu Riesza--Markowa--Kakutaniego
(por. \atlasref{ME-21::cor:measures-separate-21}{rozdzielanie miar przez funkcje ciągłe})
daje $\mu=\nu$.
\end{proof}

\begin{stwierdzenie}[Splot miar i splot miary z funkcją]
\label{prop:dt-fourier-measure-conv}
Dla skończonych zespolonych miar Radona definiujemy
$\mu*\nu$ jako obraz miary produktowej $\mu\otimes\nu$
przez odwzorowanie $(x,y)\mapsto x+y$. Wówczas
$\|\mu*\nu\|\le\|\mu\|\|\nu\|$ oraz
\[
 \widehat{\mu*\nu}=c_d^{-1}\widehat\mu\,\widehat\nu.
\]
Dla $f\in L^p(\R^d)$, $1\le p\le\infty$, splot
$\mu*f(x)=\int f(x-y)\,\mathrm d\mu(y)$ jest określony
prawie wszędzie i $\|\mu*f\|_p\le\|\mu\|\|f\|_p$.
Jeżeli $f\in L^1$, to
$\widehat{\mu*f}=c_d^{-1}\widehat\mu\,\widehat f$.
\end{stwierdzenie}
\begin{proof}
Wariacja miary produktowej jest ograniczona przez
$|\mu|\otimes|\nu|$ (można to uzyskać, przedstawiając
obie miary przez ich gęstości względem wariacji).
Wariacja obrazu nie przekracza obrazu wariacji, co
wynika z definicji przez rozłączne podziały.
Stąd oszacowanie normy splotu.
Rozkład $e^{-i(x+y)\cdot\xi}=e^{-ix\cdot\xi}e^{-iy\cdot\xi}$
i twierdzenie Fubiniego dają wzór na transformatę.

Dla $\mu\ne0$ i $1\le p<\infty$ stosujemy nierówność
Jensena względem miary probabilistycznej $|\mu|/\|\mu\|$:
\[
 \left(\int|f(x-y)|\,\mathrm d|\mu|(y)\right)^p
 \le \|\mu\|^{p-1}\int|f(x-y)|^p\,\mathrm d|\mu|(y).
\]
Po całkowaniu po $x$ i zastosowaniu Tonellego prawa strona
ma całkę $\|\mu\|^p\|f\|_p^p$. Dowodzi to również istnienia
całki prawie wszędzie. Dla $p=\infty$ wybieramy
reprezentanta ograniczonego przez $\|f\|_\infty$;
zmiana reprezentanta nie zmienia splotu prawie wszędzie
na mocy Tonellego. Dla $\mu=0$ wszystko jest oczywiste.
Jeśli $f\in L^1$, twierdzenie Fubiniego i podstawienie
$z=x-y$ dają ostatni wzór; moduł jądra podwójnej całki
ma całkę $c_d\|\mu\|\|f\|_1$.
\end{proof}

\begin{przyklad}[Miara Diraca a lemat Riemanna--Lebesgue'a]
\label{ex:dt-fourier-measure-atoms}
Mamy $\widehat{\delta_a}(\xi)=c_de^{-ia\cdot\xi}$,
więc transformata skończonej miary nie musi zanikać
w nieskończoności. Dla $\sum_{j=1}^N b_j\delta_{a_j}$
transformata jest sumą $c_d\sum_j b_je^{-ia_j\cdot\xi}$.
W przypadku gęstości całkowalnej zanikanie zapewnia
Twierdzenie~\ref{thm:dt-fourier-l1-rl}. Zatem hipoteza
o bezwzględnej ciągłości miary z całkowalną gęstością
ma tu istotne znaczenie.
\end{przyklad}
